\section{Introduction}
 
For commercial buildings, electricity costs are determined not only by
total energy consumption but also by peak power: the demand charge, levied
on the single highest power draw of a billing cycle, can account for over
30\% of the total bill [P-V2B, \S I]. In addition, utilities operate demand
response (DR) programs that provide a recurring capacity incentive to
buildings that commit to a Firm Service Level (FSL), a cap on power draw
during grid stress events, and that impose penalties when the cap is
violated during a called event. The growing presence of electric vehicles
(EVs) at workplace facilities provides, in principle, a resource suited to
both objectives: a fleet of behind-the-meter batteries that can charge
during low-price periods and discharge during peaks, reshaping the
building's net load without curtailing building services.
 
Monetizing this resource upstream, however, exposes a structural asymmetry.
The capacity payment is earned by committing, a day or more in advance, to
deliver load reduction during events whose timing the building does not
control. Unlike a stationary battery, the resource that backs this
commitment is intermittently available, since EVs arrive and depart
according to their users' schedules; self-consuming, since energy stored in
a vehicle may be expended in subsequent travel; and inter-temporally
coupled, since the state of charge (SoC) available on a peak day is
determined by charging decisions made on preceding days [P-V2B, Eq.~4].
Moreover, the resource is owned by users who are unwilling to accept
delivery commitments, penalty exposure, or prepayment schemes settled
against forecasts of their own behavior. A building that pledges
aggressively earns greater capacity revenue but bears the full delivery
risk; a building that cannot mobilize its EV fleet can pledge only the
reduction achievable through direct service curtailment, which is severely
limited for mission-critical facilities.
 
Resolving this asymmetry requires navigating a set of deeply interconnected
challenges. First, the commitment is made under forecast uncertainty: the
FSL must be pledged before building loads, EV arrivals, and, critically,
user cooperation for the day in question are realized, so the pledge
constitutes a commitment against a forecast distribution rather than a
realized schedule. Second, the deliverable resource is inter-temporally
coupled: the energy buffer available for discharge on an event day is the
residual of over-charging performed on earlier off-peak days, diminished by
off-site consumption in the interim [P-V2B, Eq.~4]. Single-day formulations
are structurally incapable of representing this channel, and the monthly
demand charge already renders rewards sparse and delayed. Third, the system
must achieve incentive alignment without user-side firmness: users must be
compensated for flexibility in departure time and required SoC, and must
not be able to profit by misreporting either [CONSENT, Thm.~1], while no
instrument available to the building may condition a user's payment or
obligation on a prediction of that user's behavior. Fourth, the resulting
delivery risk is concentrated entirely on the building: if users accept
less over-charge than forecast, or consume more of the banked energy
off-site than expected, the realized buffer is thinner than the pledge
assumed, and the penalty is borne by the operator alone. We treat this risk
not as an incidental limitation but as a quantity to be explicitly priced
and hedged.
 
Prior work addresses these challenges in isolation. Real-time V2B control
with persistence-aware buffering [P-V2B] constructs and dispatches the
inter-day buffer, but it presumes operator authority over user batteries
and earns no capacity revenue; its sole reward for banked energy is the
demand charge, a max-based functional whose marginal value saturates once
the binding monthly peak has been reduced. Strategy-proof daily negotiation
[CONSENT] prices user flexibility and guarantees voluntary participation,
but it decomposes the problem by day, so that stored energy cannot be
carried across day boundaries and the discharge depth available on a peak
day is limited to the incidental surplus with which users arrive. Two-stage
DR mechanism design [Satchidanandan et al.] formalizes the
day-ahead-commitment, real-time-delivery structure, but for resources that
are dispatchable and stationary rather than intermittently available and
self-consuming. None of these approaches, taken individually, enables a
building to assume an upstream commitment while requiring no commitment
from its users.
 
We close this gap with P-CONSENT, a three-layer architecture with a
deliberately asymmetric commitment structure. Layer~1, a day-ahead
optimization between the building and the grid, determines the FSL pledge
as its sole decision; the commitment and its associated penalty are
confined to this layer and are borne exclusively by the building. Layer~2,
a real-time Monte Carlo model predictive controller [P-V2B, Eq.~8],
delivers against the pledged FSL at minimum cost, re-optimizing as loads
and arrivals are realized. Layer~3, a negotiation mechanism [CONSENT],
generates daily user options consistent with the previously fixed FSL: on
off-peak days, the controller over-charges consenting users above their
negotiated SoC floor into the available headroom, with the excess energy
provided to the user at no cost and carried to the subsequent day [P-V2B,
Eq.~4]; on event days, it discharges the accumulated excess and compensates
users for realized peak shaving, with battery degradation incorporated into
the payment. All user-side settlement occurs on realized events only: there
is no prepayment against forecasts, no user-side commitment, and no
shortfall penalty. Within each day, information flows unidirectionally,
from the FSL to the charging controller to the negotiation mechanism; the
negotiation operates under the pledge and does not determine it. Across
days, the loop closes: the buffer accumulated on the current day, together
with the observed acceptance of over-charge offers, informs the pledge for
the following day. The persistent buffer thereby expands the FSL that the
building can safely commit, while the FSL provides the negotiation with a
concrete daily discharge target and a source of revenue, external to the
building's bill, from which user incentives are funded.
 
We state explicitly the quantity that governs the magnitude of these
benefits. The improvement attainable beyond the non-committing system is
proportional to the realizable-buffer yield: the fraction of banked excess
that survives off-site consumption and remains available for discharge,
multiplied by the fraction of users who accept the over-charge offer. Both
are behavioral quantities. Our evaluation measures them directly rather
than assuming favorable values, and we characterize how the safely
pledgeable FSL, and with it the value of the architecture, contracts as
these quantities decrease. Because the null pledge remains feasible
throughout, the system degrades gracefully to its non-committing
counterpart rather than committing an insufficient buffer to an aggressive
pledge.
 
Our contributions are:
\begin{itemize}
\item \textbf{A formulation of committed DR participation backed by an
intermittently available, self-consuming, inter-temporally coupled
resource}, with a three-layer decision structure whose within-day
information flow is unidirectional (from the FSL to the charging controller
to the negotiation mechanism) and whose cross-day loop closes through the
persistent SoC buffer (Section~\ref{sec:formulation}).
\item \textbf{A characterization of the feasible pledge set}, establishing
that the set of FSLs deliverable at a given risk level grows monotonically
with the forecast buffer, thereby formalizing the manner in which
persistence-aware over-charging expands upstream commitment capacity,
together with a pledge-depth parameter that prices the building's delivery
risk (Section~4).
\item \textbf{A realized-settlement user mechanism} that extends [CONSENT]
with cost-free over-charge offers and event-day discharge payments while
preserving budget feasibility and voluntary participation, funding improved
user offers from capacity revenue external to the building's bill; we state
precisely which of CONSENT's guarantees transfer unchanged and which
require restatement under the over-charge policy (Section~4).
\item \textbf{An evaluation on real building, tariff, and EV fleet data}
from a commercial facility (Nissan Advanced Technology Center, Silicon
Valley), measuring the realizable-buffer yield and over-charge acceptance
directly, and demonstrating that the full architecture weakly dominates
building-only, negotiation-only, and persistence-only baselines with
respect to both building and user cost, with sensitivity analyses over the
cost-sharing parameter $\alpha$ [CONSENT, Eq.~3] and the pledge depth
(Sections~6 and~7).
\end{itemize}
 
\section{Related Research}
\textit{[Deferred. Three parallel streams meeting at a gap: V2B charging
control and its daily-decomposition limit; user negotiation and incentive
mechanisms; DR mechanism design (two-stage commit/deliver; online
mechanisms for intermittent availability). Explicit non-claim on
payment-mechanism novelty.]}
 
\section{Problem Formulation}\label{sec:formulation}
 
We consider a single commercial building that manages its own EV charging
infrastructure, serves a population of persistent EV users, and
participates in a utility demand response program. The system performs
three tasks at three timescales: (i) day-ahead, it pledges a Firm Service
Level to the utility, earning a capacity payment and accepting a penalty
for violations during events; (ii) in real time, it schedules EV charging
and discharging to deliver against the pledge at minimum electricity cost;
(iii) at each user arrival, it negotiates bounded adjustments to the user's
required SoC and departure time in exchange for compensation, and offers
cost-free over-charge above the negotiated floor when headroom permits. A
notation table is provided in Table~\ref{tab:notation}; we adopt the
conventions of [P-V2B, \S III] for physical quantities and of [CONSENT,
\S 3] for negotiation objects.
 
\subsection{Specifications}\label{sec:specs}
 
\paragraph{Time structure.} The billing cycle consists of days $d \in
\Dset$. Each day is discretized into time steps $t \in \Td$ of length
$\dt$ (e.g., 15 minutes); $\Tset = \cup_d \Td$. Control decisions are taken
at the start of each step and held constant for its duration.
 
\paragraph{Persistent users and sessions.} Let $\Vset$ be the set of EV
users, each with a consistent identity across the billing cycle. Following
[P-V2B, \S III], user $v$'s visits are indexed by sessions $k \in \Kv =
\{0, \ldots, K_v\}$, where session $k$ spans $[\tarr(v,k), \tdep(v,k)]
\subseteq \Tset$ and $\Tvk$ denotes its time steps. Sessions need not occur
on consecutive days; the formulation accommodates arbitrary gaps between
visits. The SoC of EV $v$ at time $t$ is $\soc{t}{v} \in [0, \Ecap(v)]$,
expressed directly in energy units (kWh), where $\Ecap(v)$ is the battery
capacity. Within a session, the SoC evolves as
\begin{equation}\label{eq:soc}
\soc{t+1}{v} \;=\; \soc{t}{v} + \rate{t}{v} \,\dt,
\end{equation}
where $\rate{t}{v}$ (kW) is the charging rate, negative for discharge. Since at most one session is active at any time, the session index is omitted from per-step quantities. Between sessions $k$ and
$k+1$, the EV consumes $\Eext(v,k)$ energy off-site (kWh), so that
\begin{equation}\label{eq:carryover}
\Earr(v, k{+}1) \;=\; \Edep(v, k) - \Eext(v, k),
\end{equation}
which is the persistence coupling of [P-V2B, Eq.~4]; as in that
formulation, off-site usage nets out any external charging, and we assume
$\Eext(v,k) \le \Edep(v,k)$. Session-related quantities (arrival and
departure times, arrival SoC, off-site usage) are estimated from historical
data; availability is never elicited directly from users.
 
\paragraph{Requested SoC and the negotiated floor.} Upon arrival for
session $k$, each user reports a required departure SoC $\Ereq(v,k)$ (kWh)
and a departure time $\treq(v,k)$. Through the negotiation of Layer~3
(Section~4), the user may accept a revised floor $\Eneg(v,k) \le
\Ereq(v,k)$ and departure time, in exchange for compensation priced by the
user's marginal utility. The floor constitutes a
guarantee:
\begin{equation}\label{eq:floor}
\Edep(v,k) \;\ge\; \Eneg(v,k).
\end{equation}
The floor may be reduced only through the user's voluntary, compensated
agreement; it is never relaxed unilaterally.
 
\paragraph{Banked excess, survival, and the realizable buffer.} On
off-peak days, based on building load and estimated peaks realized in the problem, the controller may charge user $v$ above the negotiated floor. The banked excess at departure from session $k$ is
\begin{equation}\label{eq:banked}
\banked{v}{k} \;=\; \max\bigl(0,\; \Edep(v,k) - \Eneg(v,k)\bigr)
\quad \text{(kWh)}.
\end{equation}
Let $\Ebararr(v,k{+}1)$ denote the counterfactual arrival SoC under
$\Edep(v,k) = \Eneg(v,k)$, holding $\Eext(v,k)$ fixed; by
\eqref{eq:carryover}, $\Ebararr(v,k{+}1) = \max\bigl(0,\, \Eneg(v,k) -
\Eext(v,k)\bigr)$. The survived excess at the next arrival is
\begin{equation}\label{eq:survived}
\survived{v}{k{+}1} \;=\; \Earr(v,k{+}1) - \Ebararr(v,k{+}1)
\quad \text{(kWh)},
\end{equation}
which is nonnegative by construction, and the per-user yield is
$\yieldf{v}{k{+}1} = \survived{v}{k{+}1}/\banked{v}{k}$ whenever
$\banked{v}{k} > 0$. Each over-charge offer is accepted or declined, which
is recorded by $\accept{v}{k} \in \{0,1\}$. Let $\Vd{d} \subseteq \Vset$
denote the users with a session active on day $d$, and for $v \in \Vd{d}$
let $k$ denote that session. The realizable buffer available on day $d$ is
\begin{equation}\label{eq:buffer}
\buff{d} \;=\; \sum_{v \in \Vd{d}} \accept{v}{k{-}1}\cdot \survived{v}{k}
\quad \text{(kWh)},
\end{equation}
deliverable as power $\eta \buff{d} / (|\evwin{d}|\,\dt)$ over an event
window $\evwin{d}$, where $\eta$ is the discharge efficiency. The
population quantities $\mathbb{E}[\rho]$ and $\mathbb{E}[o]$ constitute the
central empirical question of this paper: their product, scaled by the
per-kWh value of event-window discharge, bounds the improvement achievable
by the full architecture relative to its non-committing counterpart. Both
quantities are estimated online and measured empirically in Section~7;
neither is assumed.
 
\begin{assumption}[Exogenous off-site usage]\label{as:exo}
$\Eext(v,k)$ is unaffected by the over-charge decision; that is, the
provision of cost-free energy does not induce additional vehicle usage.
This is a behavioral assumption subject to a potential rebound effect, and
we revisit it when interpreting the measured yield (Section~7).
\end{assumption}
 
\paragraph{Chargers and building load.} The building operates $\Nc$
heterogeneous chargers, bidirectional or unidirectional, with rate limits
$[\rmin{\zeta(v)}, \rmax{\zeta(v)}]$ for the charger $\zeta(v)$ assigned to $v$ ($\rmin{} < 0$ only for bidirectional chargers). Assignment follows a first-come, first-served rule and remains fixed for the duration of the session. With availability indicator $\avail{t}{v} \in \{0,1\}$, the building's total power draw is
\begin{equation}\label{eq:load}
\Ptot{t} \;=\; \bload{t} + \sum_{v \in \Vset} \avail{t}{v}\, \rate{t}{v},
\qquad \Ptot{t} \ge 0,
\end{equation}
where $\bload{t}$ (kW) is the non-EV building load.
 
\paragraph{Electricity bill.} The bill combines time-of-use energy charges
and a demand charge levied on the monthly peak:
\begin{equation}\label{eq:bill}
\ecost{t} = \we{t}\, \Ptot{t}\, \dt, \qquad
\dcost = \wdr{} \cdot \max_{t \in \Tset} \Ptot{t},
\end{equation}
with energy price $\we{t}$ (\$/kWh) and demand rate $\wdr{}$ (\$/kW). Discharging incurs a battery degradation cost $\Kbatt$ per kWh discharged, and departure below the negotiated floor incurs a penalty $\Ksoc$ per kWh.
 
\paragraph{Demand response program.} The building participates in a
capacity-paying DR program. For day $d$, it pledges a Firm Service Level
$\FSL_d$ (kW). The utility may call an event with window $\evwin{d}
\subseteq \Td$ (possibly empty); events are exogenous, sparse, and
program-dependent. The capacity payment and the violation penalty are
\begin{equation}\label{eq:dr}
\Cinc{d} = \rinc \cdot (\Phist - \FSL_d), \qquad
\Cpen{d} = \Kfsl \sum_{t \in \evwin{d}} \vio{t}\, \dt,
\end{equation}
where $\Phist$ is the program's historical baseline, $\rinc$ (\$/kW) is the
per-day-normalized capacity rate, $\vio{t} = \max(0, \Ptot{t} - \FSL_d)$ is
the violation slack, and $\Kfsl$ (\$/kWh) is the penalty rate [DR,
Eqs.~1--2]. The pledge constitutes the only commitment in the system; no
user-facing quantity is subject to commitment, penalty, or prepayment.
 
\begin{remark}[Program stylization]\label{rem:program}
Deployed programs differ in pledge cadence: the Base Interruptible Program
fixes the FSL at enrollment for a season, whereas capacity bidding programs
admit monthly nomination with day-ahead and day-of options. We model a
day-ahead re-pledgeable FSL, which corresponds most closely to the
day-ahead capacity-bidding option, and normalize the monthly capacity rate
to $\rinc$ per pledged day. Results under an enrollment-fixed FSL are a
constrained special case ($\FSL_d \equiv \FSL$ for all $d$) and are
reported in the sensitivity analysis.
\end{remark}
 
\paragraph{User cost and payments.} User $v$'s cost for session $k$ extends:
\begin{equation}\label{eq:usercost}
\ucost{k}{v} \;=\; \sum_{t \in \Tvk} \we{t}\, \gch{t}{v} \;-\;
\alpha \cdot \Umu{\theta_v} \;-\; \epay{k}{v},
\end{equation}
where $\gch{t}{v}$ is the energy charged to $v$ (negative for discharge),
$\Umu{\theta_v}$ is the user's marginal utility under the FSL-constrained
system [CONSENT, Eq.~2, evaluated under the constraint set of Layer~2],
$\alpha \in [0,1]$ is the cost-sharing parameter, and $\epay{k}{v} \ge 0$
is the event payment for the session, comprising discharge compensation and a
peak-shaving incentive computed by counterfactual replay on realized
outcomes, with the degradation cost $\Kbatt$ incorporated into the payment
so that the user does not bear battery wear separately. The banked excess
itself is provided at no cost: $\banked{v}{k}$ does not appear in
\eqref{eq:usercost}. For sessions that overlap no event window,
$\epay{k}{v} = 0$. Every term of
\eqref{eq:usercost} is settled on realized quantities; no term depends on a
forecast of user $v$'s behavior.
 
\subsection{Three-Layer Decision Structure}\label{sec:layers}
 
The system's decisions partition into three layers with distinct
timescales, information sets, and owners. A precise statement of the
information flow constitutes the structural core of the formulation.
 
\paragraph{Layer 1 (day-ahead, building--grid).} At the end of day $d-1$,
the building selects the pledge $\FSL_d$; no other decision is taken at
this layer. The selection is made with respect to the information set
\begin{equation*}
\infoDA{d} = \bigl(\text{load and arrival forecasts},\;
\text{event likelihood},\; \buffhat{d},\; \ohat,\; \rhohat\bigr),
\end{equation*}
where $\buffhat{d}$ is the forecast realizable buffer constructed from
\eqref{eq:buffer} under the current acceptance and yield estimates
$(\ohat, \rhohat)$. The pledge maximizes expected capacity revenue net of
expected penalty and delivery cost (Section~4.1). The commitment and its
penalty are confined to this layer.
 
\paragraph{Layer 2 (real-time, building control).} At each step $t$ of day
$d$, the charging controller selects the rates $\act{t} = \{\rate{t}{v}\}$
given the system state and the fixed $\FSL_d$. The controller is the
MC-MPC of [P-V2B, Eq.~8], with the pledge entering as a constraint,
$\Ptot{t} \le \FSL_d + \vio{t}$ for $t \in \evwin{d}$, penalized at rate
$\Kfsl$ (Section~4.2). The controller delivers against the pledge but does
not revise it.
 
\paragraph{Layer 3 (per-arrival, user negotiation).} At each user arrival,
the negotiation mechanism generates flexibility options and prices them by
marginal utility, with every candidate evaluated under the constraint set
of Layer~2 and hence under $\FSL_d$ (Section~4.3). On off-peak days, it
additionally extends the cost-free over-charge offer; on event days, it
schedules buffer discharge and computes $\epay{k}{v}$.
 
\paragraph{Information flow.} Within day $d$, the flow is unidirectional:
\begin{equation*}
\FSL_d \;\longrightarrow\; \text{Layer 2 constraint set}
\;\longrightarrow\; \text{Layer 3 option generation}.
\end{equation*}
The negotiation operates under the pledge and does not determine it; we
therefore make no claim of within-day bidirectional coupling, which would
render the formulation circular. Across days, the loop closes through the
buffer and the estimators:
\begin{equation*}
(\text{over-charge},\, o,\, \rho)\ \text{on day } d
\;\longrightarrow\; (\buffhat{d+1},\, \ohat,\, \rhohat)
\;\longrightarrow\; \FSL_{d+1}.
\end{equation*}
The persistence buffer expands the feasible pledge (formalized as
Proposition~1, Section~4.4), while the pledge provides the negotiation
with a concrete discharge target and an external source of revenue. This
structure, unidirectional within each day and closed across days, is what
enables the asymmetric commitment: the building can extend a firm upstream
guarantee because the pledged quantity is selected after observing the
buffer that its users have already voluntarily accumulated.
 
\paragraph{Delivery risk.} The pledge $\FSL_d$ is committed against the
estimates $(\ohat, \rhohat)$. If realized acceptance or yield falls below
its forecast, then $\buff{d} < \buffhat{d}$ and the building is exposed to
$\vio{t} > 0$ during events, with the penalty $\Kfsl$ borne entirely by the
building. The pledge-depth parameter $\lambda \in [0,1]$ (Section~4.1)
interpolates between the null pledge $\FSL_d = \Phist$, which carries zero
revenue and zero risk, and the deepest pledge feasible at risk level
$\varepsilon$. We treat $(\lambda, \varepsilon)$ as a primary operating
parameter and characterize its revenue--penalty trade-off in Section~7.
 
\subsection{Objectives}\label{sec:objectives}
 
\paragraph{Building.} The building minimizes its expected net cost over
the billing cycle:
\begin{equation}\label{eq:J}
\begin{aligned}
J = \mathbb{E}\Bigl[\,
&\sum_{t \in \Tset} \ecost{t}
+ \dcost
- \sum_{d} \Cinc{d}
+ \sum_{d} \Cpen{d} \\
&+ \Ksoc \sum_{v,k} \zund{k}{v}
+ \Kbatt \sum_{v,t} \qdis{t}{v}
+ \sum_{v,k} \bigl(\alpha\, \Umu{\theta_v} + \epay{k}{v}\bigr)
\Bigr],
\end{aligned}
\end{equation}
where $\zund{k}{v}$ denotes unmet floor energy and $\qdis{t}{v}$ denotes
discharged energy, mirroring [CONSENT, Eq.~1] augmented with the DR terms
of \eqref{eq:dr}. The expectation is taken over building load, arrivals,
events, and user responses.
 
\paragraph{Users.} Each user faces the cost \eqref{eq:usercost}. The
mechanism must preserve voluntary participation and budget feasibility
[CONSENT, Thms.~2--3], and we require that no payment, obligation, or
penalty on the user side be conditioned on a forecast: settlement occurs on
realized events only. Strategy-proofness of flexibility reporting is
inherited from [CONSENT, Thm.~1] for the negotiation step; Section~4.5
states the precise conditions under which it extends to the system with
over-charge, an extension that is not immediate.
 
\paragraph{Solution space.} A feasible solution comprises a pledge
sequence $\{\FSL_d\}$; charger assignments; per-step rates
$\{\rate{t}{v}\}$ satisfying \eqref{eq:soc}, \eqref{eq:floor},
\eqref{eq:load}, the charger limits, and the FSL constraint with slack;
negotiated agreements $\{\Eneg, \tdep, \zeta^v\}$; and event payments
$\{\epay{k}{v}\}$, jointly satisfying the participation and budget
guarantees.
